%------------------------------------------------------------------------------%
\section{Função definida por partes}

A tarifa mensal de um plano de telefonia fixa tem duas faixas de preço
\begin{itemize}
  \item Os clientes que usam até 400 minutos mensais em ligações
    pagam um valor fixo de \reais{40} por mês;

  \item Para cada minuto adicional (ou seja, que excede os 400 minutos)
    paga-se \reais{0.05}.
\end{itemize}

Se consideramos que um cliente fale $x$ minutos por mês,
podemos encontrar uma função que fornece o valor a ser pago
(por mês) da conta telefônica.

Primeiramente note que se suas ligações ultrapassarem 400 minutos, o tempo
adicional equivalerá a $x-400$.

Para determinar a função $f$ que fornece o valor mensal da conta, devemos
considerar as duas situações previstas no plano
\begin{gather*}
  \text{se } x \leq 400, \text{ então } f(x) = 40; \\
  \text{se } x    > 400, \text{ então } f(x) = 40+0,05(x-400).
\end{gather*}
Em síntese, temos a seguinte função definida por partes
\[
  f(x) = \begin{cases}
  40,            & \text{ se } x \leq 400, \\
  40+0,05(x-400) & \text{ se } x    > 400.
  \end{cases}
\]

%------------------------------------------------------------------------------%
\begin{definicao}{Função por Partes}{def:funcao-por-partes}
  Um função é dita definida por partes\index{Função!por partes}
  quando são definidas por mais de uma expressão.
\end{definicao}
%------------------------------------------------------------------------------%

%------------------------------------------------------------------------------%
\begin{example}
  Esboce o gráfico da função definida por partes
  \[
    f(x)=\begin{cases}
      2x+3, & \text{ se } x < 0,        \\
      x^2,  & \text{ se } 0 \leq x < 2, \\
      1,    & \text{ se } x \geq 2.
    \end{cases}
  \]
  {\large \color{red} Gerar o gráfico.}
\end{example}
%------------------------------------------------------------------------------%

%------------------------------------------------------------------------------%
\begin{example}
  Esboce o gráfico da função definida por partes
  \[
    g(x) = \begin{cases}
      \dfrac{x^2-4}{2-x}, & \text{ se } x \neq 2, \\
      1,                  & \text{ se } x =    2.
    \end{cases}
  \]
  {\large \color{red} Gerar o gráfico.}
\end{example}
%------------------------------------------------------------------------------%
